\section{Preliminaries}\label{sec:prel} 

\subsection{Some notation} For $n\in\N$, 
$[n]:=\{1, \dots, n\}$. The set of size-$\ell$ subsets of $[n]$ is denoted as 
$\binom{[n]}{\ell}$. 
%
%
%


\subsection{Vector spaces} For a field $\F$, $\F^n$ is the linear space 
consisting of length-$n$ 
column vectors over $\F$. We use $e_i$ to denote the $i$th standard basis vector 
in $\F^n$. 
The dual space of $\F^n$ is denoted as $(\F^n)^*$, and the dual vector 
of $v\in \F^n$ is denoted as $v^*$. 
For $S\subseteq\F^n$, $\span(S)$ 
denotes the subspace spanned by vectors in $S$. For $v\in \F^n$ and $i\in[n]$, 
$v(i)$ denotes the $i$th component of $v$. 
For $S\subseteq\F^n$, 
$S^\perp:=\{v\in \F^n \mid \forall u\in S, \tr{v}u=0\}$.


\subsection{Matrices} 
We use $\M(\ell\times n, \F)$ to denote the linear space 
of $\ell 
\times n$ 
matrices over $\F$, and set $\M(n, \F):=\M(n\times n, \F)$. For $A\in\M(\ell\times 
n, \F)$, $A(i,j)$ is the $(i,j)$th entry of $A$. We shall use $\vzero$ to denote 
all-zero 
vectors or matrices of appropriate sizes.
A matrix $A\in \M(n, 
\F)$ 
is \emph{alternating}, if for any $v\in \F^n$, $\tr{v}Av=0$. 
%
%
For 
$(i, j)\in[n]\times [n]$, $E_{i,j}\in \M(n, 
\F)$ is 
the $n\times n$ matrix with the $(i, j)$th entry being $1$, and the rest entries 
being $0$. For 
$\{i,j\}\in \binom{[n]}{2}$, $i<j$, $A_{i,j}\in 
\M(n, \F)$ is the $n\times n$ alternating matrix with the $(i,j)$th entry being 
$1$, the 
$(j,i)$th entry being $-1$, and the rest entries being $0$. 
The 
general linear group of degree $n$ over $\F$ is denoted by $\GL(n, 
\F)$. 


\subsection{Alternating matrix spaces}\label{subsec:alt_mat_sp}
Let
$\Lambda(n, \F)$ be 
the linear space of $n\times n$ alternating matrices over $\F$. Subspaces of 
$\Lambda(n, \F)$ are called alternating matrix spaces. 

%
%
%

Two alternating matrix spaces $\cA, \cB\leq \Lambda(n, \F)$ are \emph{isometric}, 
if there exists $T\in\GL(n, \F)$, such that 
$\cA=\tr{T}\cB T:=\{\tr{T}BT \mid B\in\cB\}$.

Let $\cA\leq\Lambda(n, \F)$. Suppose $U\leq \F^n$ is of dimension $d$, 
%
%
%
%
%
%
%
%
%
and let $T\in \M(n\times d, \F)$ be a matrix whose column 
vectors span $U$. The restriction of $\cA$ to $U$ via $T$ is $\cA|_{U, 
T}:=\{\tr{T}AT \mid A\in \cA\}\leq \Lambda(d, \F)$. Given another $T'\in\M(n\times 
d,\F)$ whose columns 
also 
span $U$, $\cA|_{U, T}$ and $\cA|_{U, T'}$ are isometric. Therefore, we may write 
$\cA|_U$ as the restriction of $\cA$ to $U$ via some $T\in\M(n\times d, \F)$ whose 
columns span $U$.

We say that $U$ is a \emph{totally-isotropic space}\footnote{Totally-isotropic 
spaces are also called totally singular \cite{Atk73}, isotropic \cite{BGH87}, 
and $\ell$-singular 
\cite{DS10} (for alternating $\ell$-linear maps) in the literature. We 
adopt the terminologies of totally-isotropic and 
anisotropic (see 
Section~\ref{subsec:opp_struct}), which are from the study of bilinear forms 
\cite{MH73}, and have been used in e.g. \cite{BMW17}. } for $\cA$, if 
$\dim(\cA|_U)=0$. That is, for any $u, u'\in U$ and $A\in \cA$, we have 
$\tr{u}Au'=0$. 
We say that $U$ is a \emph{complete space} for $\cA$, if 
$\dim(\cA|_U)=\binom{d}{2}$. 

Given $v\in \F^n$, the \emph{degree} of $v$ in $\cA$, $\deg_\cA(v)$, is the 
dimension of $\{ Av \mid A\in \cA\}\leq \F^n$. When $\cA$ is clear from the 
context, 
we may simply write $\deg(v)$ instead of $\deg_\cA(v)$. The \emph{minimum degree} 
of $\cA$, denoted as $\delta(\cA)$, is the minimum over the degrees over non-zero 
$v$ in $\cA$. 

Given $S\subseteq \F^n$, the \emph{radical space} of $S$ in $\cA$ is 
$\rad_\cA(S):=\{u\in \F^n \mid \forall A\in \cA, v\in S, \tr{u}Av=0\}$. We may 
simply 
write $\rad_\cA(\{v\})$ as $\rad_\cA(v)$. When $\cA$ is 
clear from the context, we may simply write $\rad(S)$ instead of $\rad_\cA(S)$. We 
also define the \emph{radical space} of $\cA\leq \Lambda(n, \F)$ as
$\rad(\cA):=\rad_\cA(\F^n)=\{v\in \F^n \mid \forall A\in\cA, Av=0\}$.

\subsection{Alternating multilinear maps} An $\ell$-linear map 
$\phi:(\F^n)^{\times \ell}\to \F^m$ is \emph{alternating}, if for any $(v_1, 
\dots, 
v_\ell)\in (\F^n)^{\times \ell}$ where $v_i=v_j$ for some $i\neq j$, $\phi(v_1, 
\dots, v_\ell)=0$.
Here, $(\F^n)^{\times \ell}$ denotes 
the $\ell$-fold Cartesian product of $\F^n$. Two alternating $\ell$-linear maps 
$\phi, \psi: (\F^n)^{\times \ell}\to \F^m$ are \emph{isomorphic}, if there exists 
$(S, 
T)\in\GL(n, \F)\times\GL(m, \F)$, such that for any $v_1, \dots, v_\ell\in \F^n$, 
$\phi(S(v_1), \dots, S(v_\ell))=T(\phi(v_1, \dots, v_\ell))$. 


%
%
%


%
%
%
%
%
%
%

\subsection{3-way arrays} Matrices are 2-way arrays, i.e. an 
array with two indices. We shall also need the notion of \emph{3-way arrays}, 
namely arrays with three indices. We use $\M(\ell\times m\times 
n, \F)$ to denote the linear space of 3-way arrays with the index set being 
$[\ell]\times[m]\times [n]$. Let $\tA\in\M(\ell\times m\times n, \F)$ be a 3-way 
array. Following \cite{KB09,GQ19}, we define the following. The \emph{frontal 
slices} of $\tA$ are $A_1, 
\dots, A_n\in \M(\ell\times m)$, where $A_k(i, j)=\tA(i,j,k)$. The \emph{tube 
fibres} of $\tA$ are $v_{i,j}\in \F^n$, $i\in[\ell]$, $j\in[m]$, where 
$v_{i,j}(k)=\tA(i,j,k)$.

3-way arrays are also referred to as 3-tensors in 
some literature. We adopt 3-way arrays, because 3-tensors are usually considered 
to be $3$-way arrays together with 
the natural action of $\GL(\ell, \F)\times\GL(m, \F)\times\GL(n, \F)$. In this 
paper, as the reader will see soon, 3-way arrays are used to record the structure 
constant of alternating 
bilinear maps. Therefore, the group action of relevance in this context is by 
$\GL(n, \F)\times 
\GL(m, \F)$ 
where $\GL(n, \F)$ acts covariantly on the first two indices. 
%
%

\subsection{Relations between 3-way arrays, bilinear maps, and matrix spaces} 
\label{subsec:rel}
It is not hard to see that alternating bilinear maps, alternating matrix spaces, 
and 3-way arrays are closely related. We spell out some details here.

An alternating bilinear map 
$\phi:\F^n\times 
\F^n\to 
\F^m$ can be represented as a tuple of alternating matrices 
$\bA=(A_1, \dots, A_m)\in\Lambda(n, \F)^m$, such that for any $u, v\in \F^n$, 
$\phi(u, v)=\tr{(\tr{u}A_1v, \dots, \tr{u}A_mv)}$.

From an alternating matrix tuple $\bA\in\Lambda(n, \F)^m$, we can construct a 
$3$-way array $\tA\in\M(n\times n\times m, \F)$ whose frontal slices are $A_i$'s. 
We can also construct an alternating matrix space $\cA=\span\{A_1, \dots, 
A_m\}\leq\Lambda(n, \F)$. 

Let $\cA\leq\Lambda(n, \F)$ be an $m$-dimensional alternating matrix space. Let
$\bA=(A_1, \dots, A_m)\in\Lambda(n, \F)$ be an ordered linear basis of $\cA$. Then 
$\bA$ gives rise to an alternating bilinear map $\phi$ and a 3-way array $\tA$ as 
above. Note that different ordered bases of $\cA$ yield different but isomorphic
alternating bilinear maps.

\section{Proof of Theorem~\ref{thm:main}}\label{sec:thm_proof}

%
%

\subsection{Restatement of Theorem~\ref{thm:main}} By fixing a basis for $U$ and 
identifying alternating bilinear forms with alternating matrices, 
Theorem~\ref{thm:main} can be restated as follows. 

\paragraph{Theorem~\ref{thm:main}, restated.} Let $\F$ be a field, $s, t\in \N$, 
$s, t\geq 2$, and $n\geq s\cdot t^4$. For any alternating matrix space
$\cA\leq\Lambda(n, \F)$, either there exists an $s$-dimensional totally-isotropic 
space of $\cA$, or there exists a $t$-dimensional complete space of $\cA$.


%
%
%
%


\subsection{Proof outline} 
To start with, by restricting to any subspace of dimension $s\cdot t^4$, we can 
assume that $\cA\leq\Lambda(n, \F)$ where $n=s\cdot t^4$. 

The proof consists of four steps. Before going into the 
details, let us first outline the objective for each step. 

%
%

\begin{description}
\item[Step 1] This step either constructs a $s$-dimensional 
totally-isotropic space of $\cA$, or computes a $n'$-dimensional $P\leq \F^n$, 
such that the minimum degree $\delta(\cA|_P)\geq t^4$. 

Let $\cB=\cA|_P\leq\Lambda(n', \F)$. The goal of the next three steps is to 
construct a dimension-$(t+1)$ complete space for $\cB$.

\item[Step 2] Let $t'=t^2$. This step
constructs a 
$Q_2\in \M(n'\times (t'+1), \F)$, of rank-$(t'+1)$, such that $\cC=\tr{Q_2}\cB 
Q_2\leq 
\Lambda(t'+1, \F)$ (the restriction of $\cB$ to the subspace of $\F^{n'}$ spanned 
by the columns of $Q_2$) contains matrices $C_1, \dots, C_{t'}\in \cC$, where
$C_i=\begin{bmatrix}
\tilde C_i & \vzero \\
\vzero & \vzero 
\end{bmatrix}$, $\tilde C_i$ has size $(i+1)\times (i+1)$, and $C_i(i, 
i+1)=1$.

\item[Step 3] Let $r=t+\binom{t}{2}$, and recall that $t'=t^2$. This step 
constructs
$Q_3\in\GL(t'+1, \F)$, such that $\cD=\tr{Q_3}\cC Q_3\leq 
\Lambda(t'+1, 
\F)$ contains 
$D_1, 
\dots, D_r\in \Lambda(t'+1, \F)$ satisfying (1) for $i\in [t]$, 
$D_i=\begin{bmatrix}
\tilde D_i & \vzero \\
\vzero & \vzero 
\end{bmatrix}$ where $\tilde D_i\in \Lambda(i+1, \F)$ and $D_i(i, i+1)=1$,  
and 
(2) for 
$i\in[\binom{t}{2}]$, $D_{t+i}=\begin{bmatrix}
\tilde D_{t+i} & \vzero & \vzero & \vzero \\
\vzero & 0 & 1 & \vzero \\
\vzero & -1 & 0 & \vzero \\
\vzero & \vzero & \vzero & \vzero 
\end{bmatrix}$, where $\tilde D_{t+i}$ is of size $t+1+2(i-1)$.

\item[Step 4] 
Based on $D_1, \dots, D_r$ from Step 3, this step constructs $Q_4\in\GL(t'+1, 
\F)$, such that $W=\span\{ e_1, \dots, e_{t+1}\}$ is a complete space for 
$\tr{Q_4}\cD Q_4$. 
It is clear that the complete space $W$ of $\tr{Q_4}\cD Q_4$ translates to a 
complete space for $\cA|_P$ through $Q_2$, $Q_3$, and $Q_4$, giving us the desired 
complete 
space for $\cA$.
%
%
%
%
%
%
%
%
%
%
%
%
%
%
%
\end{description}

Each step relies on a lemma which could be of independent interest. In the 
following, we explain these steps in detail.

\subsection{Step 1} 
%
%
The first step 
relies on the following lemma. 
Recall that the minimum degree $\delta(\cdot)$ is defined in 
Section~\ref{subsec:alt_mat_sp}.
\begin{lemma}\label{lem:iso_min-deg}
Let $s, d\in \N$, $s, d\geq 2$, and $n=s\cdot d$. For any $\cA\leq\Lambda(n, 
\F)$, either there exists a totally-isotropic space of 
dimension $s$ of $\cA$, or there exists $P\leq\F^n$, such that $\dim(P)\geq 2d$, 
and $\delta(\cA|_P)\geq 
d$. 
\end{lemma}

\begin{proof} %
Consider the following 
procedure in at most $s$ rounds. The basic 
idea is that in each round, if there exists a non-zero vector $v$ of degree $<d$, 
we restrict the current alternating matrix space to $\rad(v)$.

More specifically, in the first round, if there are no non-zero $v_1\in\F^n$ 
such that $\deg_\cA(v_1)<d$, then
$\F^n$ satisfies what we need for $P$. Otherwise, there exists a non-zero 
$v_1\in\F^n$ 
such that $\deg_\cA(v_1)<d$. Let $S_1=\span\{ v_1\}$, 
$T_1=\rad_\cA(S_1)$, and $\cA_1=\cA|_{T_1}$. By the alternating 
property, $v_1\in T_1$, so $S_1\leq T_1$. Make $R_1$ a complement subspace 
of $S_1$ in $T_1$. 
%
%
%
%
%
%
%
%
%
%
%
By $\deg_\cA(v_1)<d$, we have $\dim(T_1)\geq 
(s-1)d+1$ and $\dim(R_1)\geq (s-1)d$. Also note that $S_1\leq \rad(\cA_1)$. We 
then 
continue to 
the next round.

%
%
%
%
%
%
%
%
%
%
%
%
%
%
%
%
%
%
%
%
%

Then before the $i$th round, $i=2, \dots, s-1$, we have obtained from the 
$(i-1)$th round the following data.
\begin{itemize}
\item $S_{i-1}=\span\{ v_1, \dots, v_{i-1}\}\leq \F^n$, $\dim(S_{i-1})=i-1$, 
and $S_{i-1}$ is a totally-isotropic space of $\cA$.
\item $T_{i-1}\leq \F^n$, $\dim(T_{i-1})\geq (s-(i-1))d+(i-1)$, and $S_{i-1}\leq 
T_{i-1}$.
\item $R_{i-1}$, a complement subspace of $S_{i-1}$ in $T_{i-1}$. Note that 
$\dim(R_{i-1})\geq (s-(i-1))d\geq (s-(s-1-1))d=2d$.
\item $\cA_{i-1}=\cA|_{T_{i-1}}$, and $S_{i-1}\leq \rad(\cA_{i-1})$.
\end{itemize}
%

We then try to find a non-zero $v_i\in 
R_{i-1}$ such that 
$\deg_{\cA_{i-1}}(v_i)<d$. 
%
%
%
If no such $v_i$ exist, then $R_{i-1}$ satisfies what 
we need for $P$. 
%
Otherwise, let $S_i=\span\{ v_1, \dots, v_i\}$, 
$T_i=\rad_{\cA_{i-1}}(S_i)=\rad_{\cA_{i-1}}( v_i)$, 
and $\cA_i=\cA|_{T_i}$. Set $R_i$ to be a complement subspace of $S_i$ in $T_i$.
As $\deg_{\cA_{i-1}}(v_i)<d$, $\dim(T_i)\geq (s-i)d+i$, and $\dim(R_i)\geq 
(s-i)d$. Clearly $S_i\leq\rad(\cA_i)$, so in particular, $S_i$ is a 
totally-isotropic space for $\cA$.
%
%
We then continue 
to the 
$(i+1)$th round. 

Now suppose we just enter the $s$th round. At this point, we have 
$\dim(T_{s-1})\geq 
d+(s-1)$, and $\dim(R_{s-1})\geq d$. 
%
%
%
Take 
any non-zero $v_s\in R_{s-1}$, and set 
$S_s=\span\{ v_1, \dots, v_s\}$. Since $S_{s-1}\leq\rad(\cA_{s-1})$, 
$S_s$ is 
a totally-isotropic space of 
dimension $s$.
\end{proof}

Back to our original setting, recall that $n=s\cdot t^4$. Let $d=t^4$, so 
$n=s\cdot d$. Applying Lemma~\ref{lem:iso_min-deg} gives us 
either 
a totally-isotropic 
space of 
dimension $s$, or a subspace $P$ such that $n'=\dim(P)\geq 2d$ and 
$\delta(\cA|_P)\geq d$. In the former case, we can conclude the proof of 
Theorem~\ref{thm:main}. 
In the latter case, we will construct a 
dimension-$(t+1)$ totally-isotropic space for $\cB$ in the next three steps. 


%
%
%
%
%



\subsection{Step 2}  The second step relies on the following lemma.

\begin{lemma}\label{lem:step2}
Let $d, t'\in \N$, $d=t'^2$. Suppose $\cB\leq\Lambda(n', \F)$ satisfies that 
$n'\geq 2d$ and $\delta(\cB)\geq d$. Then there exists $Q\in \M(n'\times (t'+1), 
\F)$, of rank-$(t'+1)$, such that $\cC=\tr{Q}\cB 
Q\leq 
\Lambda(t'+1, \F)$ contains matrices $C_1, \dots, C_{t'}\in \cC$, where
$C_i=\begin{bmatrix}
\tilde C_i & \vzero \\
\vzero & \vzero 
\end{bmatrix}$, $\tilde C_i$ has size $(i+1)\times (i+1)$, and $C_i(i, 
i+1)=1$.
\end{lemma}
\begin{proof}
%
%
%
%
%
%
%
%
%
%
%
%
%

%
%
%
%
%

%
Consider the following procedure in at most $t'$ 
rounds. 

In the first round, take any non-zero $w_1\in \F^{n'}$, and find $B_1\in \cB$, 
such that $B_1w_1\neq\vzero$. Then there exists $w_2\in\F^{n'}$
such that $\tr{w_2}B_1w_1\neq 0$. Such $B_1$ exists, as $\dim(\rad_{\cB}(w_1))\leq 
n'-d$. Set $T_2=(\span\{ B_1w_1, B_1w_2\})^\perp$, and 
$W_2=\span\{ w_1, w_2\}$. By $\tr{w_2}B_1w_1\neq 0$, $T_2\cap W_2=0$. By the 
alternating property, $\dim(W_2)=2$. 

%
%
%
%
%
%
%
%
%
%
%
%
%
%
%
%
%
%
%
%
%
%
%

Then before the $i$th round, $i=2, \dots, t'$, we have obtained $w_1, 
\dots, 
w_i\in \F^{n'}$, $B_1, 
\dots, B_{i-1}\in \cB$, 
such that $\forall j\in[i-1]$, (1) $\tr{w_j}B_jw_{j+1}\neq 0$, and (2) $\forall 
1\leq 
k\leq i$, $\forall j+1<\ell\leq i$, $\tr{w_k}B_jw_\ell=0$ . We 
also have $T_i=(\span\{ B_jw_k \mid j\in[i-1], k\in[i]\})^\perp$, and $W_i=\span\{ 
w_1, \dots, w_i\}$, such that $T_i\cap W_i=0$. 

We claim that there exist 
$w_{i+1}\in T_i$ and $B_i\in \cB$, such that $\tr{w_{i+1}}B_iw_i\neq 0$. To see 
this, note that $\dim(T_i)\geq n'-(i-1)\cdot i$ and $\dim(\rad_\cB(w_i))\leq 
n'-d$. So as long as $d>(i-1)\cdot i$, we can take any $w_{i+1}\in T_i\setminus 
\rad_{\cB}(w_i)$, for which there exists $B_i\in \cB$ such that 
$\tr{w_{i+1}}B_iw_i\neq 0$. 

Let $T_{i+1}=(\span\{ B_jw_k \mid j\in[i], k\in[i+1]\})^\perp$, and 
$W_{i+1}=\span\{ w_1, \dots, w_{i+1}\}$. We claim that $T_{i+1}\cap 
W_{i+1}=0$. If not, suppose 
$w=\alpha_1w_1+\dots+\alpha_iw_i+\alpha_{i+1}w_{i+1}\in 
T_{i+1}$. Let $j$ be the smallest integer such that $\alpha_j\neq 0$. If $j\leq 
i$, then $\tr{w}B_jw_{j+1}$ is non-zero. If $j=i+1$, then $\tr{w}B_iw_i$ is 
non-zero. In 
either case, this is a contradiction to the assumption that $w\in T_{i+1}$. We 
also note that, by examining $\tr{w_{i+1}}B_iw_i$, we have $B_i\not\in\span\{ B_1, 
\dots, B_{i-1}\}$. 

We perform the above operations, and after the $t'$-th round we get the desired 
$w_1, 
\dots, w_{t'+1}$ and $B_1, \dots, B_{t'}\in\Lambda(n', \F)$. This requires 
$d>(t'-1)t'$, which is 
fine as we have set $d=t'^2$. 

Let $Q=\begin{bmatrix} w_1 & \dots & 
w_{t'+1}\end{bmatrix}\in\M(n'\times(t'+1),\F)$, 
and let $\cC=\tr{Q}\cB Q\leq\Lambda(t'+1, \F)$. We claim that $\cC$ satisfies the 
requirements of this lemma. Recall that for any $i\in[t']$, $B_i$ satisfies that 
$\tr{w_i}B_iw_{i+1}\neq 0$. Furthermore, for any $i+1<\ell\leq t+1$ and $1\leq 
k\leq 
t+1$, $\tr{w_k}B_iw_\ell=0$. Set $\tr{w_i}B_iw_{i+1}=\alpha_i$. Let
$C_i=\tr{Q}(\frac{1}{\alpha_i} B_i)Q\in \cC$. Then $C_i(i, 
i+1)=\frac{1}{\alpha_i}\tr{w_i}B_iw_{i+1}=1$, and for any $i+1<\ell\leq t+1$ and 
$1\leq k\leq t+1$, $C_i(k, \ell)=\frac{1}{\alpha_i}\tr{w_k}B_iw_\ell=0$. 
Intuitively, this just 
means 
that $C_i$ is of the form $\begin{bmatrix}
\tilde C_i & \vzero\\
\vzero & \vzero
\end{bmatrix}$, where $\tilde C_i$ is of size $(i+1)\times (i+1)$, and $C_i(i, 
i+1)=1$. 
\end{proof}

From Step 1, we have $\cB\leq\Lambda(n', \F)$ with 
$n'\geq 
2d$ and $\delta(\cB)\geq d$. Recall that $d=t^4$, and set $t'=t^2$. Applying 
Lemma~\ref{lem:step2} to $\cB$, $d$, $t'$ produces $Q_2\in\M(n'\times(t'+1), \F)$ 
which achieves the objective of this step. 

\subsection{Step 3} The objective of this step is fulfilled by the following 
lemma. 
\begin{lemma}\label{lem:step3}
Let $t'=t^2$, and $r=t+\binom{t}{2}$. 
Suppose $\cC\leq 
\Lambda(t'+1, \F)$ contains matrices $C_1, \dots, C_{t'}\in \cC$, where
$C_i=\begin{bmatrix}
\tilde C_i & \vzero \\
\vzero & \vzero 
\end{bmatrix}$, $\tilde C_i$ has size $(i+1)\times (i+1)$, and $C_i(i, 
i+1)=1$.
Then there exists $Q\in\GL(t'+1, \F)$, such that $\cD=\tr{Q}\cC Q\leq 
\Lambda(t'+1, 
\F)$ contains 
$D_1, 
\dots, D_r\in \Lambda(t'+1, \F)$ satisfying (1) for $i\in [t]$, 
$D_i=\begin{bmatrix}
\tilde D_i & \vzero \\
\vzero & \vzero 
\end{bmatrix}$ where $\tilde D_i\in \Lambda(i+1, \F)$ and $D_i(i, i+1)=1$, 
and 
(2) for 
$i\in[\binom{t}{2}]$, $D_{t+i}$ is of the form
%
%
%
%
%
%
\begin{equation}\label{eq:form_of_D}
\begin{bmatrix}
\tilde D_{t+i} & \vzero & \vzero & \vzero & \dots & \vzero \\
\vzero & 0 & 1 & 0 & \dots & 0 \\
\vzero & -1 & 0 & 0 & \dots & 0 \\
\vzero & 0 & 0 & 0 & \dots & 0 \\
\vdots & \vdots & \vdots & \vdots & \ddots & \vdots \\
\vzero & 0 & 0 & 0 & \dots & 0 
\end{bmatrix}, 
\end{equation}
where $\tilde D_{t+i}\in\Lambda(t+1+2(i-1), \F)$. 
\end{lemma}

That is, in the matrix $D_{t+i}$, the only nonzero 
entries 
of the $(t+2i)$th and $(t+2i+1)$th rows and columns are at the $(t+2i, t+2i+1)$ 
and $(t+2i+1, t+2i)$ positions. The reason for imposing this condition will be 
clear later from Observation~\ref{obs:key}.


\begin{proof}[Proof of Lemma~\ref{lem:step3}]
Observe that the $(t+2i, t+2i+1)$ and 
$(t+2i+1, t+2i)$ entries of $C_{t+2i}$ are $1$ and $-1$, respectively. So we can 
put $C_{t+2i}$ in the desired form, by multiplying appropriate elementary 
matrices  
(i.e. $I+\alpha\cdot E_{t+2i, j}$ for $j<t+2i$ and appropriate $\alpha\in\F$) on 
the left and their transposes on the right, 
to set other entries on the $(t+2i+1)$th row and 
column to be $0$. 
%
Some 
care is required to ensure that during 
this process, other $C_{t+2j}$'s, if they are already in this form, are not 
affected. 

Therefore, we apply appropriate matrices to 
%
$C_{t+2i}$, 
for $i$ in a \emph{decreasing} order, namely starting from $i=\binom{t}{2}$ and 
then 
going to $i=1$. To see that this does not affect those $C_{t+2j}$'s which were 
already in this form, let us examine $C_{t+2i}$, $C_{t+2i+1}$, and $C_{t+2i+2}$, 
when we put 
$C_{t+2i}$ into the form as in (\ref{eq:form_of_D}). Note that $C_{t+2i+2}$ 
has 
been put into the desired form as in (\ref{eq:form_of_D}). That is, 
$${\small
C_{t+2i}=\begin{bmatrix}
\hat C_{t+2i} & \vstar & \vstar & \vzero & \vzero& \dots & \vzero \\
\vstar & 0 & 1 & 0 & 0 & \dots & 0 \\
\vstar & -1 & 0 & 0 & 0 & \dots & 0 \\
\vzero & 0 & 0 & 0 & 0 & \dots & 0 \\
\vzero & 0 & 0 & 0 & 0 & \dots & 0 \\
\vdots & \vdots & \vdots & \vdots &\vdots & \ddots & \vdots \\
\vzero & 0 & 0 & 0 & 0 &\dots & 0 
\end{bmatrix},}$$
$${\small
C_{t+2i+1}=\begin{bmatrix}
\hat C_{t+2i+1} & \vstar & \vstar &  \vstar & \vzero & \dots & \vzero \\
\vstar & 0 & * & * & 0 &\dots & 0 \\
\vstar & * & 0 & 1 & 0 & \dots & 0 \\
\vstar & * & -1 & 0 & 0 & \dots & 0 \\
\vzero & 0 & 0 & 0 & 0 & \dots & 0 \\
\vdots & \vdots & \vdots & \vdots &\vdots & \ddots & \vdots \\
\vzero & 0 & 0 & 0 & 0 &\dots & 0 
\end{bmatrix},}$$
$${\small 
C_{t+2i+2}=\begin{bmatrix}
\hat C_{t+2i+2} & \vstar & \vstar & \vzero & \vzero & \dots & \vzero \\
\vstar & 0 & * & 0 & 0 & \dots & 0 \\
\vstar & * & 0 & 0 & 0 &\dots & 0 \\
\vzero & 0 & 0 & 0 & 1 &\dots & 0 \\
\vzero & 0 & 0 & -1 & 0 &\dots & 0 \\
\vdots & \vdots & \vdots & \vdots & \vdots &\ddots & \vdots \\
\vzero & 0 & 0 & 0 & 0 &\dots & 0 
\end{bmatrix},}
$$
where $\hat C_{t+2i}$, $\hat C_{t+2i+1}$, and $\hat C_{t+2i+2}$ are in 
$\Lambda(s+2i-1, \F)$.
From the above, when we use $(t+2i, t+2i+1)$ and $(t+2i+1, t+2i)$ entries  
to set other entries on the $(t+2i)$th and $(t+2i+1)$th rows and columns in 
$C_{t+2i}$ to be zero, such operations
do not affect the $(t+2i+2)$th and $(t+2i+3)$th rows and columns of 
$C_{t+2i+2}$, nor the $(t+2j)$th and $(t+2j+1)$th rows and columns of $C_{t+2j}$ 
for $j\geq i+1$ in general. Furthermore, such operations do not change $C_{t+2j}$ 
for $j\leq i-1$ at all. It follows that after these operations, all $C_{t+2i}$, 
$i\in[\binom{t}{2}]$, are in the form of (\ref{eq:form_of_D}) as desired.

%
%
%
%
%
%
%

Suppose $(C_1, \dots, C_{t'})$ are changed to $(C_1', 
\dots, 
C_{t'}')$ after these operations, which implicitly define $Q\in\GL(t'+1, \F)$. We 
then do the following. For $i\in[t]$, let $D_i=C_i'$. For 
$i\in[\binom{t}{2}]$, let $D_{t+i}=C_{t+2i}'$. We then obtain 
$r=t+\binom{t}{2}=\binom{t+1}{2}$ matrices 
$D_1, \dots, D_r\in\Lambda(t'+1, \F)$ with $t'=t^2$ which satisfy the properties 
as required by this lemma, concluding the proof.
\end{proof}

%
%
%
%
%
%
%
%
%
%
%
%

\subsection{Step 4} To start with, we need an observation on complete 
spaces as follows. 
%
%
Given an $m$-dimensional $\cA\leq\Lambda(n, \F)$, let
$\bA=(A_1, 
\dots, A_m)\in \Lambda(n, \F)^m$ be an ordered basis of $\cA$. 
From $\bA$, we construct a $3$-way 
array $\tA\in\M(n\times n\times m, \F)$ whose frontal slices are $A_i$'s. Let 
$f_{i,j}\in \F^m$, $i, j\in[n]$, be the tube fibres of $\tA$. Let
$W=\span\{ e_1, \dots, e_t\}\leq \F^n$, where $e_i$'s are standard basis 
vectors. We then note the following characterisation of $W$ to be a complete space 
of $\cA$.
\begin{observation}\label{obs:tube}
Let $\cA$, $f_{i,j}$, and $W$ be as above. Then $W$ is a complete space of 
$\cA$, 
if and only 
if, $f_{i,j}$, 
$1\leq i<j\leq t$, are linearly independent. 
\end{observation}

The following lemma completes Step 4. 
\begin{lemma}\label{lem:step4}
Let $t'=t^2$, and $r=t+\binom{t}{2}$.
Suppose $\cD\leq \Lambda(t'+1, 
\F)$ contains 
$D_1, 
\dots, D_r\in \Lambda(t'+1, \F)$, such that (1) for $i\in [t]$, 
$D_i=\begin{bmatrix}
\tilde D_i & \vzero \\
\vzero & \vzero 
\end{bmatrix}$ where $\tilde D_i\in \Lambda(i+1, \F)$ and $D_i(i, i+1)=1$, 
and 
(2) for 
$i\in[\binom{t}{2}]$, $D_{t+i}$ is of the form
%
%
%
%
%
%
$i\in[\binom{t}{2}]$, $D_{t+i}=\begin{bmatrix}
\tilde D_{t+i} & \vzero & \vzero & \vzero \\
\vzero & 0 & 1 & \vzero \\
\vzero & -1 & 0 & \vzero \\
\vzero & \vzero & \vzero & \vzero 
\end{bmatrix}$,
where $\tilde D_{t+i}\in\Lambda(t+1+2(i-1), \F)$. Then there exists $Q\in\GL(t'+1, 
\F)$, such that $W=\span\{ e_1, \dots, e_{t+1}\}$ is a complete space for 
$\tr{Q}\cD Q$. 
\end{lemma}

%
%
%
%
%
%
\begin{proof}
Construct a $3$-way array $\tD$ of size $(t'+1)\times(t'+1)\times r$, where the 
$i$th frontal slice is $D_i$. Let $f_{i,j}\in\F^r$ be the $(i, j)$th tube fibre of 
$\tD$. Note that the tube fibres $f_{1, 2}, f_{2,3}, \dots, f_{t, t+1}$ and 
$f_{t+2, t+3}, f_{t+4, t+5}, \dots, f_{t', t'+1}$ are linearly 
independent. 


Let us arrange the tube fibres $f_{i,j}$, $1\leq i\leq t-1$, $i+2\leq j\leq t$, in 
the reverse lexicographic order, and relabel them accordingly as $\tilde f_k$ for 
$k\in[\binom{t}{2}]$. That is, $\tilde f_1=f_{1,3}$, $\tilde f_2=f_{1,4}$, $\tilde 
f_3=f_{2, 4}$, $\tilde f_4=f_{1, 5}$, $\tilde f_5=f_{2,5}$, and so on.

%
%
Our goal is to apply appropriate elementary matrices (on the left and their 
transposes on the right) to $\tD$ to make the tube 
fibres at positions $(k, \ell)$, $1\leq k<\ell\leq t+1$, linearly independent.
If this could 
be achieved, Observation~\ref{obs:tube} ensures that $W=\span\{ e_1, \dots,
e_{t+1}\}$ is a complete space of the resulting alternating matrix space.

To do that, consider the following operations in $\binom{t}{2}$ rounds. After the 
$i$th round, we wish to maintain that 
$$f_{1, 2}, f_{2,3}, \dots, f_{t, t+1}, \tilde f_1, 
\dots, \tilde f_i, f_{t+2(i+1), t+2i+3}, f_{t+2(i+2), t+2i+5}, \dots, f_{t', 
t'+1}$$
are linearly 
independent. Note that $t'=t+2\cdot \binom{t}{2}$. If this could be achieved, 
after $\binom{t}{2}$ rounds, 
$$f_{1,2}, f_{2,3}, \dots, f_{t, t+1}, \tilde f_1, \dots, \tilde 
f_{\binom{t}{2}}$$ 
would be linearly independent.

Recall that before the first round starts, we have that the tube fibres $f_{1, 
2}$, 
$f_{2,3}$, $\dots$, $f_{t, t+1}$ and 
$f_{t+2, t+3}, f_{t+4, t+5}, \dots, f_{t', t'+1}$ are linearly 
independent. 

Suppose now we have completed the $i$th round. Let us explain the operations in 
the 
$(i+1)$th round. We first check if 
$$f_{1, 
2}, f_{2,3}, \dots, f_{t, t+1}, \tilde f_1, 
\dots, \tilde f_i, \tilde f_{i+1}, f_{s+2(i+2), s+2i+5}, 
f_{s+2(i+3), s+2i+7}, \dots, f_{t', 
t'+1},$$ 
are linearly independent. 

If so, we proceed to the next round. 

If not, we have that $\tilde f_{i+1}$ is in the linear span of 
$$f_{1, 
2}, f_{2,3}, \dots, f_{t, t+1}, \tilde f_1, 
\dots, \tilde f_i, f_{t+2(i+2), t+2i+5}, 
f_{t+2(i+3), t+2i+7}, \dots, f_{t', 
t'+1}.$$ 
Now we 
wish to add $f_{t+2(i+1), t+2i+3}$ to 
$\tilde f_{i+1}$. This is because, since after the $i$th round we had that
\begin{equation}\label{eq:cond}
f_{1, 2}, f_{2,3}, 
\dots, f_{t, t+1}, \tilde f_1, 
\dots, \tilde f_i, f_{t+2(i+1), t+2i+3}, f_{t+2(i+2), t+2i+5}, \dots, f_{t', 
t'+1},
\end{equation}
are linearly independent, we have that 
\begin{equation}\label{eq:desired}
f_{1, 2}, f_{2,3}, 
\dots, f_{t, t+1}, \tilde f_1, 
\dots, \tilde f_i, \tilde f_{i+1}+f_{t+2(i+1), t+2i+3}, f_{t+2(i+2), t+2i+5}, 
\dots, f_{t', 
t'+1},
\end{equation}
are also linearly independent. 

But we cannot add $f_{t+2(i+1), t+2i+3}$ to 
$\tilde f_{i+1}$ directly. Indeed, the legitimate operations are left multiplying  
elementary matrices (namely $I+E_{i,j}$) and right multiplying 
their transposes. 
These 
correspond to performing row and column operations on $\tD$ viewed as a matrix 
$(f_{i,j})_{i,j\in[t'+1]}$ whose entries are vectors. We will make use of this 
perspective in the following. 

Suppose $\tilde f_{i+1}$ corresponds to 
$f_{j, k}$ for 
some $1\leq j<k\leq t+1$. In order to add $f_{t+2(i+1), t+2i+3}$ to 
$\tilde f_{i+1}$, we can first add the $(t+2(i+1))$th row to the $j$th row, and to 
maintain the alternating property, add 
the $t+2(i+1)$th column to the $j$th column. Then we add the $(t+(2i+3))$th column 
to the $k$th column, and to maintain the alternating property, add the 
$(t+(2i+3))$th row to the $k$th row. This does add $f_{t+2(i+1), t+2i+3}$ to 
$\tilde f_{i+1}$. 

However, it is 
possible that during the above procedure, some of the vectors in 
$$\{f_{1, 
2}, f_{2,3}, \dots, f_{t, t+1}, \tilde f_1, 
\dots, \tilde f_i\}$$ get altered as well. 
(It is easy to see that $f_{t+2(i+2), t+2i+5}, \dots, f_{t', t'+1}$ are not 
changed.)
For example, when adding the $(t+2(i+1))$th row to the $j$th row, $f_{j, j+1}$ and 
those $\tilde f_{i'}$ corresponding to $f_{j, j'}$ could be added by certain 
vectors as 
well. 
Therefore, instead of getting those vectors in (\ref{eq:desired}), we 
 get
\begin{multline}\label{eq:actual}
 f_{1, 2}+g_{1,2}, f_{2,3}+g_{2,3}, 
\dots, f_{t, t+1}+g_{t,t+1}, \tilde f_1+\tilde g_1, 
\dots, \tilde f_i+\tilde g_i, \\
 \tilde f_{i+1}+f_{t+2(i+1), t+2i+3}+\tilde g_{i+1}, f_{t+2(i+2), t+2i+5}, \dots, 
 f_{t', t'+1},
\end{multline}
where $g_{j,j+1}$ and $\tilde g_k\in \F^r$. 

Therefore, we need to show that those 
vectors in (\ref{eq:actual}) are linearly independent. The following 
observation is crucial for this. 
\begin{observation}\label{obs:key}
We have that $g_{j, j+1}$ and $\tilde 
g_k$ are in $\span\{f_{t+2(i+2), t+2i+5}$, $\dots$, $f_{t', t'+1}\}$. 
\end{observation}
\begin{proof}
Note that $g_{j, j+1}$ and $\tilde g_k$ come from those fibre tubes $f_{p, 
t+2(i+1)}$ and $f_{q, t+2i+3}$, where $1\leq p, q\leq t+1$. As can be seen from 
(\ref{eq:form_of_D}), 
these fibre tubes are in the linear span of $f_{t+2(i+2), t+2i+5}, \dots, f_{t', 
t'+1}$, because the only non-zero entries on the $(t+2i+2)$th and $(t+2i+3)$th 
columns of $D_{t+i+1}$ are in the $(t+2i+3, t+2i+2)$ and $(t+2i+2, t+2i+3)$th 
positions.
\end{proof}

Given 
this observation, the linear independence of vectors in (\ref{eq:actual})
follows from the linear independence of vectors in (\ref{eq:cond}). Indeed, 
by a change of basis, we can assume that the vectors in (\ref{eq:cond}) 
form a set of standard basis vectors in the order they are listed. So putting 
those vectors in (\ref{eq:cond}) as 
column vectors in  
a matrix form simply gives
$$
\begin{bmatrix}
I_{t+i} & \vzero &\vzero \\
\vzero & 1 & \vzero \\
\vzero & \vzero & I_{\binom{t+1}{2}-(t+i+1)}
\end{bmatrix},
$$
where $I_k$ denotes the $k\times k$ identity 
matrix. Now putting the vectors in (\ref{eq:actual}) as column vectors in a 
matrix gives
$$
\begin{bmatrix}
I_{t+i} & {\bf *} &\vzero \\
\vzero & 1 & \vzero \\
\mathbf{*} & * & I_{\binom{t+1}{2}-(t+i+1)}
\end{bmatrix},
$$
where $\mathbf{*}$ means that the entries could be arbitrary. This matrix is 
clearly of full-rank, proving the linear independence of vectors in 
(\ref{eq:actual}). Note that the 
entries in the lower-left submatrix comes from $g_{i, i+1}$ and $\tilde g_j$, and 
the entries in the $(t+i+1)$-th column comes from $f_{t+2(i+1), t+(2i+3)}+\tilde 
g_{i+1}$. This 
also explains the necessity of Step 3, as otherwise the 
upper-left $(t+i+1)\times(t+i+1)$ submatrix would be of the form
$\begin{bmatrix}
I_{t+i} & {\bf *}\\
{\bf *} & 1
\end{bmatrix}$, which could be not full-rank. 


Now that we achieved what we wanted in the $(i+1)$th round, also note that the 
above operations do not 
affect the $2j$ and $2j+1$ rows and columns of $D_{t+j}$ for $j>i+1$. This means 
that we can have the same set up to perform the above operations (with increased 
row and column indices) in the next round. This concludes the proof of 
Lemma~\ref{lem:step4}.
\end{proof}

This 
concludes the proof of Theorem~\ref{thm:main}.\hfill\qed
